\documentclass{amsart}
% For dealing with references we use the comment environment
\usepackage{verbatim}
\newenvironment{reference}{\comment}{\endcomment}
%\newenvironment{reference}{}{}
\newenvironment{slogan}{\comment}{\endcomment}
\newenvironment{history}{\comment}{\endcomment}

% For commutative diagrams we use Xy-pic
\usepackage[all]{xy}

% We use 2cell for 2-commutative diagrams.
\xyoption{2cell}
\UseAllTwocells

% We use multicol for the list of chapters between chapters
\usepackage{multicol}

% This is generally recommended for better output
\usepackage{lmodern}
\usepackage[T1]{fontenc}

% For cross-file-references

% Package for hypertext links:
\usepackage{hyperref}

% For any local file, say "hello.tex" you want to link to please

% Theorem environments.
%
\theoremstyle{plain}
\newtheorem{theorem}[subsection]{Theorem}
\newtheorem{proposition}[subsection]{Proposition}
\newtheorem{lemma}[subsection]{Lemma}

\theoremstyle{definition}
\newtheorem{definition}[subsection]{Definition}
\newtheorem{example}[subsection]{Example}
\newtheorem{exercise}[subsection]{Exercise}
\newtheorem{situation}[subsection]{Situation}

\theoremstyle{remark}
\newtheorem{remark}[subsection]{Remark}
\newtheorem{remarks}[subsection]{Remarks}

\numberwithin{equation}{subsection}

% Macros
%
\def\lim{\mathop{\mathrm{lim}}\nolimits}
\def\colim{\mathop{\mathrm{colim}}\nolimits}
\def\Spec{\mathop{\mathrm{Spec}}}
\def\Hom{\mathop{\mathrm{Hom}}\nolimits}
\def\Ext{\mathop{\mathrm{Ext}}\nolimits}
\def\SheafHom{\mathop{\mathcal{H}\!\mathit{om}}\nolimits}
\def\SheafExt{\mathop{\mathcal{E}\!\mathit{xt}}\nolimits}
\def\Sch{\mathit{Sch}}
\def\Mor{\mathop{\mathrm{Mor}}\nolimits}
\def\Ob{\mathop{\mathrm{Ob}}\nolimits}
\def\Sh{\mathop{\mathit{Sh}}\nolimits}
\def\NL{\mathop{N\!L}\nolimits}
\def\CH{\mathop{\mathrm{CH}}\nolimits}
\def\proetale{{pro\text{-}\acute{e}tale}}
\def\etale{{\acute{e}tale}}
\def\QCoh{\mathit{QCoh}}
\def\Ker{\mathop{\mathrm{Ker}}}
\def\Im{\mathop{\mathrm{Im}}}
\def\Coker{\mathop{\mathrm{Coker}}}
\def\Coim{\mathop{\mathrm{Coim}}}

% Boxtimes
%
\DeclareMathSymbol{\boxtimes}{\mathbin}{AMSa}{"02}

%
% Macros for moduli stacks/spaces
%
\def\QCohstack{\mathcal{QC}\!\mathit{oh}}
\def\Cohstack{\mathcal{C}\!\mathit{oh}}
\def\Spacesstack{\mathcal{S}\!\mathit{paces}}
\def\Quotfunctor{\mathrm{Quot}}
\def\Hilbfunctor{\mathrm{Hilb}}
\def\Curvesstack{\mathcal{C}\!\mathit{urves}}
\def\Polarizedstack{\mathcal{P}\!\mathit{olarized}}
\def\Complexesstack{\mathcal{C}\!\mathit{omplexes}}
% \Pic is the operator that assigns to X its picard group, usage \Pic(X)
% \Picardstack_{X/B} denotes the Picard stack of X over B
% \Picardfunctor_{X/B} denotes the Picard functor of X over B
\def\Pic{\mathop{\mathrm{Pic}}\nolimits}
\def\Picardstack{\mathcal{P}\!\mathit{ic}}
\def\Picardfunctor{\mathrm{Pic}}
\def\Deformationcategory{\mathcal{D}\!\mathit{ef}}

\setlength{\emergencystretch}{3em}
\begin{document}
\title[Stacks Project, Tag 06H0]{396 --- Essential Artinian surjections characterized by cotangent differentials}
\maketitle
\noindent Copyright (C) 2005--2025 Johan de Jong. Stacks Project authors. Source: \href{https://stacks.math.columbia.edu/tag/06H0}{Tag 06H0}.

\noindent Editorial difficulty note (not author text). Difficulty H2: The converse constructs the kernel of a projected derivation as a subalgebra and proves it still surjects onto the target. Essential surjectivity then forces the projected differential kernel to vanish. The final small-extension argument turns nonessentiality into an actual section through a length comparison. .

\noindent Editorial provenance note (not author text). History: excerpt from revision \texttt{a0444\allowbreak 6e57e\allowbreak c1fbc\allowbreak 252a8\allowbreak 71afc\allowbreak ec775\allowbreak 2fb28\allowbreak 07b14}, formal-defos.tex, lines 669--781. Reference presentation was adapted; original-author context and core support blocks were retained or added. The mathematical wording is unchanged. Licensed under GNU FDL 1.2 or later; no invariant sections or cover texts. See ../licenses/COPYING. Context blocks reproduce author definitions and supporting results; they are not additional counted problems.

\begin{definition}
\label{definition-CLambda}
Let $\Lambda$ be a Noetherian ring and let $\Lambda \to k$ be a finite
ring map where $k$ is a field. We define {\it $\mathcal{C}_\Lambda$} to be
the category with
\begin{enumerate}
\item objects are pairs $(A, \varphi)$ where $A$ is an Artinian local
$\Lambda$-algebra and where $\varphi : A/\mathfrak m_A \to k$ is a
$\Lambda$-algebra isomorphism, and
\item morphisms $f : (B, \psi) \to (A, \varphi)$ are local $\Lambda$-algebra
homomorphisms such that $\varphi \circ (f \bmod \mathfrak m) = \psi$.
\end{enumerate}
We say we are in the {\it classical case} if $\Lambda$ is a Noetherian
complete local ring and $k$ is its residue field.
\end{definition}

\section{Notation and Conventions}
\label{section-notations-conventions}

\noindent
A ring is commutative with $1$. The  maximal ideal of a local ring $A$
is denoted by $\mathfrak{m}_A$. The set of positive integers is denoted
by $\mathbf{N} = \{1, 2, 3, \ldots\}$. If $U$ is an object of a
category $\mathcal{C}$, we denote by $\underline{U}$
the functor
$\Mor_\mathcal{C}(U, -): \mathcal{C} \to \textit{Sets}$, see
Remarks \href{https://stacks.math.columbia.edu/tag/06GK}{remarks cofibered groupoids (Tag 06GK)} (\href{https://stacks.math.columbia.edu/tag/06GQ}{item definition yoneda (Tag 06GQ)}).
Warning: this may conflict with the notation in other chapters where we
sometimes use $\underline{U}$ to denote $h_U(-) = \Mor_\mathcal{C}(-, U)$.

\medskip\noindent
Throughout this chapter $\Lambda$ is a Noetherian ring and
$\Lambda \to k$ is a finite ring map from $\Lambda$ to a field.
The kernel of this map is denoted $\mathfrak m_\Lambda$ and the
image $k' \subset k$. It turns out that $\mathfrak m_\Lambda$ is
a maximal ideal, $k' = \Lambda/\mathfrak m_\Lambda$ is a field, and
the extension $k/k'$ is finite. See discussion surrounding
(\href{https://stacks.math.columbia.edu/tag/06S2}{equation k prime (Tag 06S2)}).




\begin{definition}
\label{definition-essential-surjection}
Let $f: B \to A$ be a ring map in $\mathcal{C}_\Lambda$.  We say $f$
is an {\it essential surjection} if it has the following properties:
\begin{enumerate}
\item $f$ is surjective.
\item If $g: C \to B$ is a ring map in $\mathcal{C}_\Lambda$ such that
$f \circ g$ is surjective, then $g$ is surjective.
\end{enumerate}
\end{definition}

\begin{definition}
\label{definition-small-extension}
Let $f: B \to A$ be a ring map in $\mathcal{C}_\Lambda$.  We say $f$
is a {\it small extension} if it is surjective and $\Ker(f)$ is a nonzero
principal ideal which is annihilated by $\mathfrak{m}_B$.
\end{definition}

\noindent
To analyze essential surjections in $\mathcal{C}_\Lambda$ a bit more
we introduce some notation. Suppose that $A$ is an object
of $\mathcal{C}_\Lambda$ or more generally any $\Lambda$-algebra
equipped with a $\Lambda$-algebra surjection $A \to k$.
There is a canonical exact sequence
\begin{equation}
\label{equation-sequence}
\mathfrak m_A/\mathfrak m_A^2 \xrightarrow{\text{d}_A}
\Omega_{A/\Lambda} \otimes_A k \to
\Omega_{k/\Lambda} \to 0
\end{equation}
see
Algebra, Lemma \href{https://stacks.math.columbia.edu/tag/00RU}{lemma differential seq (Tag 00RU)}.
Note that $\Omega_{k/\Lambda} = \Omega_{k/k'}$ with $k'$ as
in (\href{https://stacks.math.columbia.edu/tag/06S2}{equation k prime (Tag 06S2)}). Let $H_1(L_{k/\Lambda})$
be the first homology module of the naive cotangent complex of $k$
over $\Lambda$, see
Algebra, Definition \ref{algebra-definition-naive-cotangent-complex}.
Then we can extend (\ref{equation-sequence})
to the exact sequence
\begin{equation}
\label{equation-sequence-extended}
H_1(L_{k/\Lambda}) \to
\mathfrak m_A/\mathfrak m_A^2 \xrightarrow{\text{d}_A}
\Omega_{A/\Lambda} \otimes_A k \to
\Omega_{k/\Lambda} \to 0,
\end{equation}
see
Algebra, Lemma \href{https://stacks.math.columbia.edu/tag/00S2}{lemma exact sequence NL (Tag 00S2)}.
If $B \to A$ is a ring map in $\mathcal{C}_\Lambda$
or more generally a map of $\Lambda$-algebras equipped
with $\Lambda$-algebra surjections onto $k$, then we obtain a
commutative diagram
\begin{equation}
\label{equation-sequence-functorial}
\vcenter{
\xymatrix{
H_1(L_{k/\Lambda}) \ar[r] \ar@{=}[d] &
\mathfrak m_B/\mathfrak m_B^2 \ar[r]_{\text{d}_B} \ar[d] &
\Omega_{B/\Lambda} \otimes_B k \ar[r] \ar[d] &
\Omega_{k/\Lambda} \ar[r] \ar@{=}[d] & 0 \\
H_1(L_{k/\Lambda}) \ar[r] &
\mathfrak m_A/\mathfrak m_A^2 \ar[r]^{\text{d}_A} &
\Omega_{A/\Lambda} \otimes_A k \ar[r] &
\Omega_{k/\Lambda} \ar[r] & 0
}
}
\end{equation}
with exact rows.

\begin{definition}
\label{algebra-definition-naive-cotangent-complex}
Let $R \to S$ be a ring map. The {\it naive cotangent complex}
$\NL_{S/R}$ is the chain complex (\href{https://stacks.math.columbia.edu/tag/07BM}{equation naive cotangent complex (Tag 07BM)})
$$
\NL_{S/R} = \left(I/I^2 \longrightarrow \Omega_{R[S]/R} \otimes_{R[S]} S\right)
$$
with $I/I^2$ placed in (homological) degree $1$ and
$\Omega_{R[S]/R} \otimes_{R[S]} S$ placed in degree $0$. We will denote
$H_1(L_{S/R}) = H_1(\NL_{S/R})$\footnote{This module is sometimes
denoted $\Gamma_{S/R}$ in the literature.} the homology in degree $1$.
\end{definition}

\begin{lemma}
\label{lemma-essential-surjection}
Let $f: B \to A$ be a ring map in $\mathcal{C}_\Lambda$.
Notation as in (\ref{equation-sequence-functorial}).
\begin{enumerate}
\item The equivalent conditions of
Lemma \href{https://stacks.math.columbia.edu/tag/06S3}{lemma surjective (Tag 06S3)}
characterizing when $f$ is surjective are also equivalent to
\begin{enumerate}
\item $\Im(\text{d}_B) \to \Im(\text{d}_A)$ is surjective, and
\item the map $\Omega_{B/\Lambda} \otimes_B k \to
\Omega_{A/\Lambda} \otimes_A k$ is surjective.
\end{enumerate}
\item The following are equivalent
\begin{enumerate}
\item $f$ is an essential surjection
(see Lemma \href{https://stacks.math.columbia.edu/tag/06S5}{lemma essential surjection mod squares (Tag 06S5)}),
\item the map $\Im(\text{d}_B) \to \Im(\text{d}_A)$ is an
isomorphism, and
\item the map $\Omega_{B/\Lambda} \otimes_B k \to
\Omega_{A/\Lambda} \otimes_A k$ is an isomorphism.
\end{enumerate}
\item If $k/k'$ is separable, then $f$ is an essential surjection if
and only if the map
$\mathfrak m_B/(\mathfrak m_\Lambda B + \mathfrak m_B^2) \to
\mathfrak m_A/(\mathfrak m_\Lambda A + \mathfrak m_A^2)$
is an isomorphism.
\item If $f$ is a small extension, then $f$ is not essential if and only if
$f$ has a section $s: A \to B$ in $\mathcal{C}_\Lambda$
with $f \circ s = \text{id}_A$.
\end{enumerate}
\end{lemma}

\begin{proof}
Proof of (1). It follows from (\ref{equation-sequence-functorial})
that (1)(a) and (1)(b) are equivalent. Also, if
$A \to B$ is surjective, then (1)(a) and (1)(b) hold. Assume (1)(a).
Since the kernel of $\text{d}_A$ is the image of
$H_1(L_{k/\Lambda})$ which also maps to
$\mathfrak m_B/\mathfrak m_B^2$ we conclude that
$\mathfrak m_B/\mathfrak m_B^2 \to \mathfrak m_A/\mathfrak m_A^2$
is surjective. Hence $B \to A$ is surjective by
Lemma \href{https://stacks.math.columbia.edu/tag/06S3}{lemma surjective (Tag 06S3)}. This finishes the proof of (1).

\medskip\noindent
Proof of (2). The equivalence of (2)(b) and (2)(c) is immediate from
(\ref{equation-sequence-functorial}).

\medskip\noindent
Assume (2)(b). Let $g : C \to B$ be a ring map in $\mathcal{C}_\Lambda$
such that $f \circ g$ is surjective. We conclude that
$\mathfrak m_C/\mathfrak m_C^2 \to \mathfrak m_A/\mathfrak m_A^2$
is surjective by
Lemma \href{https://stacks.math.columbia.edu/tag/06S3}{lemma surjective (Tag 06S3)}.
Hence
$\Im(\text{d}_C) \to \Im(\text{d}_A)$ is surjective
and by the assumption we see that
$\Im(\text{d}_C) \to \Im(\text{d}_B)$ is surjective.
It follows that $C \to B$ is surjective by (1).

\medskip\noindent
Assume (2)(a). Then $f$ is surjective and we see that
$\Omega_{B/\Lambda} \otimes_B k \to \Omega_{A/\Lambda} \otimes_A k$
is surjective. Let $K$ be the kernel. Note that
$K = \text{d}_B(\Ker(\mathfrak m_B/\mathfrak m_B^2 \to
\mathfrak m_A/\mathfrak m_A^2))$ by (\ref{equation-sequence-functorial}).
Choose a splitting
$$
\Omega_{B/\Lambda} \otimes_B k =
\Omega_{A/\Lambda} \otimes_A k \oplus K
$$
of $k$-vector space. The map $\text{d} : B \to \Omega_{B/\Lambda}$
induces via the projection onto $K$ a map $D : B \to K$. Set
$C = \{b \in B \mid D(b) = 0\}$. The Leibniz rule shows that this is
a $\Lambda$-subalgebra of $B$. Let $\overline{x} \in k$. Choose $x \in B$
mapping to $\overline{x}$. If $D(x) \not = 0$, then we can find an element
$y \in \mathfrak m_B$ such that $D(y) = D(x)$. Hence $x - y \in C$ is
an element which maps to $\overline{x}$. Thus $C \to k$ is surjective
and $C$ is an object of $\mathcal{C}_\Lambda$. Similarly, pick
$\omega \in \Im(\text{d}_A)$. We can find $x \in \mathfrak m_B$
such that $\text{d}_B(x)$ maps to $\omega$ by (1). If $D(x) \not = 0$, then
we can find an element $y \in \mathfrak m_B$ which maps to zero
in $\mathfrak m_A/\mathfrak m_A^2$ such that $D(y) = D(x)$.
Hence $z = x - y$ is an element of $\mathfrak m_C$ whose
image $\text{d}_C(z) \in \Omega_{C/k} \otimes_C k$ maps to $\omega$.
Hence $\Im(\text{d}_C) \to \Im(\text{d}_A)$ is surjective.
We conclude that $C \to A$ is surjective by (1). Hence $C \to B$ is
surjective by assumption. Hence $D = 0$, i.e., $K = 0$, i.e., (2)(c) holds.
This finishes the proof of (2).

\medskip\noindent
Proof of (3). If $k'/k$ is separable, then
$H_1(L_{k/\Lambda}) =
\mathfrak m_\Lambda/\mathfrak m_\Lambda^2 \otimes_{k'} k$, see
Lemma \href{https://stacks.math.columbia.edu/tag/06S9}{lemma H1 separable case (Tag 06S9)}.
Hence $\Im(\text{d}_A) =
\mathfrak m_A/(\mathfrak m_\Lambda A + \mathfrak m_A^2)$
and similarly for $B$. Thus (3) follows from (2).

\medskip\noindent
Proof of (4). A section $s$ of $f$ is not surjective (by definition a
small extension has nontrivial kernel), hence $f$ is not essentially
surjective. Conversely, assume $f$ is a small extension but not an
essential surjection. Choose a ring map $C \to B$ in $\mathcal{C}_\Lambda$
which is not surjective, such that $C \to A$ is surjective. Let
$C' \subset B$ be the image of $C \to B$. Then $C' \not = B$ but
$C'$ surjects onto $A$. Since $f : B \to A$ is a small extension,
$\text{length}_C(B) = \text{length}_C(A) + 1$. Thus
$\text{length}_C(C') \leq \text{length}_C(A)$ since
$C'$ is a proper subring of $B$. But $C' \to A$ is surjective, so in
fact we must have $\text{length}_C(C') = \text{length}_C(A)$ and
$C' \to A$ is an isomorphism which gives us our section.
\end{proof}
\end{document}
